\documentclass[11pt]{article}
\usepackage[margin=1in]{geometry}
\usepackage{amsmath,amssymb,hyperref}
\title{Transferable issue rights across settlement rejoins}
\date{24 November 2026}
\begin{document}
\maketitle

A warehouse group uses a support until the receiving rack physically accepts its load. An identity-bound disposition record can be issued only after that acceptance. The successor gate must answer exact issuance and learned settlement queries when separated groups exchange records and rejoin. The \href{https://github.com/eternal-november/eos-research-programmes/pull/352}{single-issuer frontier} does not immediately extend: an allocator can physically settle A:1 before granting B:1 and never exceed budget one, but separated gates receiving only the grants rejoin with width two. Physical settlement is not evidence at the gate.

\section{Issuance and visible facts}
Fix stable issuer namespaces, each with a unique consecutive counter and no reset. There are $B$ named single-spender rights, also indexing distinct physical supports; one issue spends one right. A handoff disables local spending at send and enables it at a matching receive. It carries both right identity and the prior-use vector. It can take place while old work is live and creates neither clearance nor permission to spend the right again. Authentication, one-use handoff, and physical acceptance are trusted primitives, not cryptographic claims. All receipts, grants and joins are from one non-forking history. The authority sets no global grant order: a right is ordered by ownership/spend events and a namespace by its consecutive assignments. Incomparable grants on independent rights/issuers can proceed independently.

The owner of right $r$ keeps $L_r(i)$, the largest prior issue index paid for by $r$ in namespace $i$, or zero. Spending requires locally learned certified settled prefixes covering this vector. A grant at $(i,n)$ publishes the componentwise maximum $P(i,n)$ of the pre-spend vectors for all grants through $n$ in issuer $i$. Each reported prefix is justified by actual settled facts already learned at its relevant spend; only the newest such bound per namespace is needed. The cumulative payload is stronger evidence than an ordinary header-only grant.

A gate stores the issued upper bound $U_i$, largest contiguously learned settled index $F_i$, and settled exceptions above $F_i$. A grant raises a bound atomically with its certificate. A standalone prefix may raise issuance and settlement together. A genuine individual receipt no higher than known $U_i$ can become visible immediately. If it exceeds $U_i$, immediate visibility requires its covering grant and dependent facts in one publication, or a separately charged sparse cache. Returning unknown while such a receipt is buffered preserves only a weaker visibility interface. Awaiting means issued but not learned settled, not necessarily physically occupied.

\section{Exact merge and width}
Disclosure means the issuance prefixes and the learned certified disposition facts, not a snapshot of all plant state. In each namespace take the maximum of $U_i$ and $F_i$, unite above-frontier exceptions, and normalize repeatedly through contiguous learned facts. This yields exactly union of disclosed facts, independently of reassociation, replay order or duplication. It does not justify exchanging a permit arrival with a query or an acceptance with a conflicting physical start.

Call an ID counted if $F_i<j\le U_i$, even if its receipt is a settled exception. A counted ID must be exposed by an arrived grant whose issuer prefix reaches it: an authenticated settled prefix raising $U_i$ also raises $F_i$, and a stand-alone admitted individual receipt cannot raise $U_i$. If two counted IDs share a right, let $a$ be earlier and $b$ later on that right. The grant exposing $b$ conveys $b$'s pre-spend vector within its issuer-cumulative certificate, hence a prefix covering $a$. This contradicts $a$ counted. The counted IDs inject into the $B$ rights, proving
\[
                            \sum_i(U_i-F_i)\le B
\]
for all visible views and their joins. With fixed counters and $s$ nonempty windows, the first hole in each window is known and a further $B-s$ bitmap bits suffice. This is not total bounded memory: counters, proofs, buffers and authority state are charged separately. With $k$ namespaces, active pre-spend vectors have at most $kB$ entries and issuer accumulators at most $k^2$ entries, excluding evidence representation.

On one right, reuse through A:1, B:1, C:1 after successive physical settlements. Send only the first and last grants. C:1's immediate-predecessor-only evidence covers B:1 but not A:1, giving width two. For a different control, spend $p$ on X:1 and later A:1, then independently spend $q$ on A:2. Receive only X:1 and A:2. The latter reveals A:1 by prefix, but lineage only for $q$ cannot clear X:1. Width becomes three rather than two. Both failures are reached with real physical acceptance before reuse. Bare out-of-prefix receipt disclosure can reproduce the second failure even with sound bundled grants.

\paragraph{Exact condition for per-right uniqueness.}
Restrict to immutable grant atoms exposing their issued prefixes and carried certified settlement, with their payload fixed at issue. Per-right uniqueness over every grant-subset view is equivalent to this condition: for every grant atom $g$ and each ID $b$ it exposes that it does not cover as settled, $g$ must cover every earlier spend of $b$'s right. Sufficiency follows from the injection proof. If this condition fails for predecessor $a$, deliver only $g$ and the atom issued at $a$. The latter proves neither $a$ settled (its payload is fixed at issue), nor $b$ settled (later on that same right), while $g$ covers neither. Then both IDs are counted. This proves necessity for per-right uniqueness, not minimal metadata or necessity for the weaker total bound with unused rights.

\section{Costs and direct comparators}
One missing first receipt can stop all $B$ rights. Let one issuer spend them on A:1 through A:$B$ and learn physical settlements of A:2 through A:$B$ but not A:1. Every right's last spend is above the known frontier. All are unavailable to issue even though $B-1$ physical supports are already empty. If the first receipt arrives $h$ ticks after the others, issuance is blocked throughout that interval; upon its arrival all become eligible. Thus arbitrary delays prevent a finite service bound, and the $h$-tick wait is attained for this protocol. Allowing reuse on a settled exception changes the guarantee: join the old A header across its missing first hole with the new namespace grant reusing the same right. The former use still falls within the counted window despite a receipt for it. Alternative encodings are not excluded.

Fixed escrow with permanent namespace allotments $B_i$ has total join width bounded by $\sum_i B_i$ by the original per-issuer proof. In the replay one issuer needs two simultaneous starts after another issuer's work settles; one right each refuses one, whereas local handoff enables both with no new support. Replayable transfer itself can lag indefinitely. There is no uniform dominance over fixed escrow.

A factored direct comparator uses identical local spending, physical support and certified prefix facts, and publishes the grant only after its separate dependency proof is present. On matched visible cuts its exact join and query agree with the bundle. For header/proof arrival at $a,b$, visibility occurs at $\max(a,b)$, giving wait $(b-a)^+$ relative to header-only visibility. In the three-grant example, separate units including an explicit empty proof marker use six deliveries with two nonzero vector entries; bundled units use three deliveries with the same entries. The direct baseline can omit empty markers when absence is certified or coalesce the messages, tying the bundle. The replay checks peak three unpaired parts. This accounting excludes actual signatures, retries, proof replication and network bandwidth; it is not an optimality claim. Equivalent atomic causal delivery discharges the same proof obligations.

\section{Checks and boundaries}
The replay exercises in-flight right transfer, old-owner cutoff, stale handoff rejection, false receipt rejection, physical-before-documentary acceptance, both missing dependencies and atomic receipt repair. It compares 128 split grant-subset joins with learned set union and exact per-ID queries, executes 4,320 paired header/proof cuts and 50 recurrent splits and rejoins (100 issues and 100 handoffs). Trace assertions check physical occupancy and unique authority. The general width and delay bounds above are written arguments, not conclusions from the finite checks.

\href{https://arxiv.org/abs/1503.09052}{Bounded Counter CRDTs} provide a direct precedent for transferable rights, and \href{https://arxiv.org/abs/1603.01529}{causal delta merge} for dependency-aware publication. The project's \href{https://github.com/eternal-november/eos-research-programmes/blob/e1348964f73704212d2d8b48d5982a5ed600a2ff/projects/contextual-composition/foundations/native-whole-scope-feedback/Nonblocking_Recurrent_Assembly_Reading_Route.tex}{nonblocking recurrent assembly} motivates the distinct commitment, support and read cuts; the componentwise merge here is only of compatible documentary views. No comparative proof-work advantage, liveness under arbitrary partitions, practical authentication claim or full C2 acceptance is established.
\end{document}
